\clearpage
\textbf{ЗАКЛЮЧЕНИЕ}

В рамках данной курсовой работы была проведена комплексная оценка
физических генераторов случайных последовательностей (ФГСП), основанных
на различных источниках энтропии. Изучение теоретических основ позволило
выделить важные свойства случайных чисел, отличать истинную случайность
от псевдослучайности, а также понять роль энтропии как ключевого ресурса
в криптографических и инженерных системах. Особое внимание было уделено
методам статистической проверки качества случайных последовательностей,
включая тесты частоты, последовательности бит, оценки энтропии и другие
критерии из NIST STS.

Во второй главе были подробно рассмотрены физические источники энтропии,
применяемые в аппаратных ГСП. Электронные источники включали тепловой,
лавинный, дробовой и мерцающий шумы. Для каждого из них рассмотрены
принципы генерации, аппаратная реализация, устойчивость к внешним
факторам и особенности практического применения. Установлено, что такие
источники, как лавинный и дробовой шум, позволяют реализовывать
компактные и сравнительно простые устройства, однако при этом требуют
учета технических ограничений, таких как чувствительность к температуре
или необходимость в усилении сигнала. Тепловой шум обеспечивает хорошую
воспроизводимость и низкую стоимость, в то время как мерцающий шум слабо
пригоден к использованию из-за сложности выделения полезного сигнала.

Оптические и квантовые источники, описанные во второй части главы,
открывают перспективы для создания генераторов с высокой энтропийной
мощностью и принципиально непредсказуемым поведением. Их реализация, как
правило, требует сложной оптоэлектронной аппаратуры, но зато
обеспечивает высокую теоретическую безопасность и потенциальную
устойчивость к широкому классу атак. Отдельно отмечены возможности
использования эффектов квантовой интерференции, туннелирования и
спонтанной эмиссии фотонов в качестве надёжных первоисточников
случайности.

Анализ механических и акустических источников показал, что, несмотря на
их простоту и низкую стоимость, данные методы обладают низкой
стабильностью, значительной зависимостью от внешней среды и невысокой
скоростью генерации. Однако в специфических условиях (например, в
автономных системах или устройствах сбора случайных данных из окружающей
среды) такие источники могут найти применение.

Глава, посвящённая постобработке, подчеркнула важность цифровой
обработки для повышения качества получаемых последовательностей.
Указано, что даже физически случайные сигналы могут содержать
коррелированные участки, требующие использования хеш-функций,
декоррелирующих фильтров и нормализаторов. Методы оценки и устранения
ошибок критически важны при внедрении ФГСП в ответственных
криптографических системах.

Сравнительный анализ, представленный в третьей главе, обобщил полученные
данные, выделив достоинства и недостатки каждого типа источников. По
совокупности характеристик наилучшие показатели стабильности и скорости
продемонстрировали лавинные и квантовые генераторы. В то же время
дробовой шум оказался сбалансированным решением по совокупности
критериев: надёжности, простоты реализации и чувствительности к внешним
воздействиям. Также была обозначена высокая чувствительность некоторых
источников к помехам, что требует проведения экранирования или
компенсации шумов при проектировании.

Наконец, в разделе о перспективах развития рассмотрены технологические и
теоретические направления, способные повлиять на развитие аппаратных
генераторов в ближайшие годы: это миниатюризация компонентов, внедрение
в стандартные микросхемы, повышение энергоэффективности и
самотестируемости, а также дальнейшее развитие квантовых источников как
наиболее надёжной формы получения энтропии.

Таким образом, проведённое исследование подтвердило актуальность и
сложность выбора подходящего источника энтропии для проектирования
надёжных аппаратных генераторов. Полученные результаты могут служить
основой для разработки и оптимизации ФГСП, адаптированных под конкретные
задачи в области информационной безопасности, криптографии и инженерных
систем.

